\documentclass{article}
\usepackage{amsmath,amssymb,amsfonts,amsthm}
\usepackage{multicol}
\usepackage{multirow}
\usepackage{mathtools}
\usepackage{soul}
\usepackage{hyperref}
\hypersetup{
    colorlinks=true,
    linkcolor=blue,
    filecolor=magenta,      
    urlcolor=cyan,
    pdftitle={Overleaf Example},
    pdfpagemode=FullScreen,
    }
\usepackage{color}
\usepackage[table]{xcolor}
\usepackage[T1]{fontenc}
\usepackage{etoolbox}
\usepackage{multicol}
\usepackage{multirow}
\usepackage{fancyhdr}
\usepackage{graphicx}
\usepackage{array}
\usepackage{amsthm}
\usepackage{titlesec}
\usepackage{tikz}
\usetikzlibrary{arrows.meta,calc}
\renewcommand{\baselinestretch}{1.2}

\titleformat*{\section}{\large\bfseries}
\titleformat*{\subsection}{\normalsize\bfseries}

\graphicspath{{C:/Users/teoso/OneDrive/Documents/Tugas Kuliah/Template Math Depart/}}

\theoremstyle{definition}
\newtheorem{theorem}{Theorem}
\newtheorem*{teorema}{Teorema}
\newtheorem*{definisi}{Definisi}
\newtheorem*{bukti}{Bukti}

\newcommand{\Arg}{\text{Arg}}
\fancyhead[L]{\textit{Teosofi Hidayah Agung}}
\fancyhead[R]{\textit{5002221132}}
\pagestyle{fancy}

\begin{document}
\begin{definisi}[Fungsi Cantor]
  Definisikan barisan fungsi $C_n:[0,1]\to[0,1]$ secara rekursif sebagai berikut. Ambil
  \[
    C_0(x)=x.
  \]
  Untuk $n\ge0$, definisikan
  \[
    C_{n+1}(x)=
    \begin{cases}
      \frac12 C_n(3x),           & 0\le x\le \frac13,       \\
      \frac12,                   & \frac13\le x\le \frac23, \\
      \frac12+\frac12 C_n(3x-2), & \frac23\le x\le1.
    \end{cases}
  \]
  Fungsi Cantor $C:[0,1]\to[0,1]$ didefinisikan sebagai
  \[
    C(x)=\lim_{n\to\infty}C_n(x).
  \]
\end{definisi}

\begin{teorema}
  Fungsi Cantor $C:[0,1]\to[0,1]$ adalah measurable.
\end{teorema}

\begin{bukti}
  Misalkan $\Sigma=\mathcal T=\sigma(\{(a,b)\subset[0,1]\})$.
  Menurut definisi, fungsi $f:(X,\Sigma)\to(Y,\mathcal T)$
  measurable jika $f^{-1}(E)\in\Sigma$ untuk setiap $E\in\mathcal T$. Karena $\mathcal T$ dihasilkan oleh interval terbuka,
  cukup dibuktikan bahwa $C^{-1}((a,b))\in\Sigma$. Ambil interval $(a,b)\subset[0,1]$ dan definisikan
  \[
    A=C^{-1}((a,b))=\{x\in[0,1]:a<C(x)<b\}.
  \]
  Karena fungsi Cantor monoton tidak menurun,
  maka himpunan $\{x:C(x)<b\}$ dan $\{x:C(x)\le a\}$
  masing-masing merupakan interval.
  Akibatnya $A$ juga merupakan interval $(\alpha,\beta)$
  untuk suatu $\alpha,\beta$. Karena interval berada dalam $\Sigma$,
  maka $C^{-1}((a,b))\in\Sigma$.
\end{bukti}
\end{document}